Put \(N=N_m\), \(S=\{(g_m(q),f(q)):q\in K_{\epsilon}\}\), and \(H=\operatorname{conv}(S)\subset\mathbb R^{N+1}\).  The set \(S\) is compact because \(K_{\epsilon}\) is compact and \((g_m,f)\) is continuous.  Every point of \(H\) can be reduced to a convex combination of at most \(N+1\) points of \(S\): all points of \(S\) lie in the affine hyperplane \(\sum_c x_c=1\), of dimension \(N\), and if a combination has more than \(N+1\) positive weights, an affine dependence among its support points permits changing the weights in a nonzero dependence direction until one weight vanishes, without changing the represented point.  Iteration proves the claim.  Therefore \(H\) is the continuous image of the compact set obtained from \(K_{\epsilon}^{N+1}\) and the \((N+1)\)-simplex of weights, so \(H\) is compact.

The attainable moment--target pairs are exactly \(H\): every finite convex combination is induced by its atomic probability measure.  Conversely, if the barycenter of \((g_m,f)\) under a probability measure were outside the closed convex set \(H\), Euclidean projection on \(H\) would give an affine functional strictly separating that barycenter from every point of \(S\); integrating the strict inequality would be a contradiction.  Thus
\[
 H=\left\{\left(\int g_m\,d\nu,\int f\,d\nu\right):
          \nu\in\mathcal P(K_{\epsilon})\right\}.
\]
For feasible \(b\), the section \(H_b=\{t:(b,t)\in H\}\) is nonempty, compact, and convex.  Hence it is the attained interval \([L_m(b),U_m(b)]\).  The converse construction in lem:scm-moment-reduction shows that every point of this interval, and in particular each endpoint, is attained by a compatible hierarchical SCM; thus the interval is sharp for the causal target rather than merely for an observed-law relaxation.

Now suppose \(b\in\operatorname{relint}(\mathcal B_{m,\epsilon})\).  Work in \(\operatorname{aff}(\mathcal B_{m,\epsilon})\times\mathbb R\), thereby discarding normals that merely encode \(\sum_c x_c=1\).  The point \((b,L_m(b))\) is a lower boundary point of the compact convex set \(H\).  A supporting affine functional exists directly as follows.  Project \((b,L_m(b)-k^{-1})\notin H\) onto \(H\), normalize the vector from the exterior point to its projection, and take a convergent subsequence.  The projection inequality passes to the limit and produces a nonzero pair \((a,\beta)\) such that
\[
 a^{\top}(x-b)+\beta\{t-L_m(b)\}\geq0
 \quad\text{for every }(x,t)\in H.
\]
Its last coefficient is nonnegative by the direction of the exterior sequence.  It cannot be zero: if \(\beta=0\), then \(a\) supports \(\mathcal B_{m,\epsilon}\) at its relative-interior point \(b\), which forces \(a=0\) in the affine tangent space, contrary to normalization.  Rescale so that \(\beta=1\), and define
\[
 \lambda=-a+\{L_m(b)+a^{\top}b\}\mathbf 1.
\]
Since \(\mathbf1^{\top}g_m(q)=\sum_cB_c^{(m)}(q)=1\), evaluation of the support inequality at \((g_m(q),f(q))\) gives
\[
 \lambda^{\top}g_m(q)=L_m(b)-a^{\top}\{g_m(q)-b\}\leq f(q),
 \qquad \lambda^{\top}b=L_m(b).
\]
Weak duality gives \(\lambda'^{\top}b\leq L_m(b)\) for every feasible minorant \(\lambda'\), so this \(\lambda\) attains the lower dual optimum.  Applying the same argument to the upper boundary point, with exterior points above it, supplies an attained majorant with value \(U_m(b)\).  This proves \(D_m(b)=(L_m(b),U_m(b))\) with no gap.

For the lower endpoint, let \(T=\{q:\lambda^{\top}g_m(q)=f(q)\}\).  If \(\nu\) attains the endpoint, then
\[
 0=\int\{f(q)-\lambda^{\top}g_m(q)\}\,d\nu(q),
\]
whose integrand is continuous and nonnegative; hence \(\nu(T)=1\) and \(b\in\operatorname{conv}g_m(T)\).  If the affine dimension of \(g_m(T)\) is \(r\), the same affine-dependence elimination used above reduces a representation of \(b\) to at most \(r+1\) contact points.  Their atomic measure remains feasible and has objective \(\lambda^{\top}b=L_m(b)\).  Since \(g_m(T)\subset\{x:\sum_cx_c=1\}\), \(r+1\leq N\).  The upper endpoint is identical with the majorant contact set.  The SCM reconstruction in lem:scm-moment-reduction turns both atomic measures into compatible endpoint models, proving all support assertions.